\documentclass[10pt,a4paper]{amsart}
\usepackage[margin=19mm]{geometry}
\usepackage{amsmath,amssymb,mathtools}
\usepackage[scr=rsfs,cal=euler]{mathalfa}
\usepackage{xfrac}
\usepackage[unicode,hidelinks,bookmarks=false]{hyperref}
\newcommand{\ie}{i.e.\ }
\newcommand{\cf}{cf.\ }
\newcommand{\resp}{respectively\ }
\newcommand{\Subsection}{\subsection}
\providecommand{\nobreakdash}{\nobreak\hbox{-}}
\setlength{\emergencystretch}{2em}
\multlinegap0pt
\usepackage{bm}
\usepackage{bbm}
\usepackage{dsfont}
\usepackage{xspace}
\usepackage{cancel}
\usepackage{mathtools}
\usepackage{cleveref}
\usepackage{amsthm}
\makeatletter
\@ifundefined{theorem}{\newtheorem{theorem}{Theorem}}{}
\@ifundefined{definition}{\newtheorem{definition}[theorem]{Definition}}{}
\makeatother
\DeclareMathOperator{\Frac}{Frac}
\DeclareMathOperator{\Proj}{Proj}
\DeclareMathOperator{\Spec}{Spec}
\newcommand{\Loj}{\mathcal{L}} % Lojasiewicz exponent symbol
\newcommand{\QQ}{\mathbb{Q}}
\newcommand{\RR}{\mathbb{R}}
\newcommand{\R}{\mathcal{R}}
\newcommand{\abul}{{\a_\bullet}}
\newcommand{\bbul}{{\b_\bullet}}
\newcommand{\m}{\mathfrak m}
\newcommand{\ord}{\operatorname{\mathrm{ord}}}
\renewcommand{\a}{\mathfrak{a}}
\renewcommand{\b}{\mathfrak{b}}
\title[An algebraic theory of Lojasiewicz exponents]{An algebraic theory of Lojasiewicz exponents}
\author{Tai Huy Ha}
\date{Source: arXiv:2602.18410v2 (CC BY 4.0)}
\begin{document}
\maketitle

\noindent\emph{Source and licence.} An algebraic theory of Lojasiewicz exponents. Version arXiv:2602.18410v2 (CC BY 4.0), distributed under the Creative Commons Attribution 4.0 International licence, as recorded in the arXiv licence metadata for this exact version and shown on the abstract page https://arxiv.org/abs/2602.18410v2. Full rights notice: licenses/arxiv-2602.18410-CC-BY-4.0.txt.\par\smallskip

\noindent\textit{Editorial scope.} Counted unit: the Theorem whose source label is \texttt{thm:finite-testing-fingen} in the article (source0.tex lines 1055--1073, source label \texttt{thm:finite-testing-fingen}), delivered together with the article's own complete original proof of it (source0.tex lines 1075--1200, one \texttt{proof} environment of 126 lines). No mathematics is written, repaired or re-derived: statement and proof are the article's own text line for line. Required context retained uncounted: def:L-families (lines 490--502); thm:valuative-filtrations (lines 904--912). Every definition, display and label of those lines is retained unchanged. Omitted from this package and not redistributed: 1--489, 503--903, 913--1054, 1074--1074, 1201--5162. Nothing inside a retained excerpt is omitted or altered: every retained body is delivered exactly as the article's lines are written, so no omission receipt of any kind accompanies this package. Citations: the retained excerpts carry no citation command at all, so no bibliography entry of the article is transcribed and no reference is redistributed. Reference adaptation: none was needed and none was made; every reference inside the delivered unit resolves inside it, so no unresolved reference or citation is produced. Printed slips kept literal: no printed word, symbol, hypothesis or number of the counted statement or of its proof was altered. Layout and class: the article's own class is \texttt{amsart} with packages including amsmath, amsthm, amssymb, amsfonts, mathrsfs, mathtools, thmtools, enumitem, xcolor, hyperref, cleveref, fullpage; the wrapper is the lane's pinned \texttt{amsart} editorial wrapper at 10pt on A4, which restates the theorem-like environments the retained text needs and adds the article's own macro definitions that the retained text needs (\texttt{\textbackslash Frac}, \texttt{\textbackslash Loj}, \texttt{\textbackslash Proj}, \texttt{\textbackslash QQ}, \texttt{\textbackslash R}, \texttt{\textbackslash RR}, \texttt{\textbackslash Spec}, \texttt{\textbackslash a}, \texttt{\textbackslash abul}, \texttt{\textbackslash b}, \texttt{\textbackslash bbul}, \texttt{\textbackslash m}, \texttt{\textbackslash ord}), with extra wrapper packages (bm, bbm, dsfont, xspace, cancel, mathtools, cleveref). The delivered numbering is the unit's own. Assets: no retained span contains a figure environment, a graphics-inclusion command, a PDF-page inclusion, a table or a TikZ picture; no image file is copied into this package, the package declares no asset dependency and no image file is redistributed with it. Licence: version arXiv:2602.18410v2 of this article is distributed under the Creative Commons Attribution 4.0 International licence, as recorded in the arXiv licence metadata for this exact version and shown on the abstract page https://arxiv.org/abs/2602.18410v2. A full rights notice is delivered at licenses/arxiv-2602.18410-CC-BY-4.0.txt. Characters: every retained excerpt is pure ASCII, so the delivered engine, XeLaTeX with its default Latin Modern text font, reports no missing character.
\begin{definition}[Algebraic \L ojasiewicz exponent for filtrations]
	\label{def:L-families}
	Let $\abul=\{\a_p\}_{p\ge1}$ and $\bbul=\{\b_q\}_{q\ge1}$ be filtrations
	of ideals in $R$.
	The \emph{algebraic \L ojasiewicz exponent of $\abul$ with respect to $\bbul$} is
	\[
	\Loj_{\bbul}(\abul)
	\;:=\;
	\inf\left\{\frac{q}{p}\in\QQ_{>0}\ :\ \b_{qt}\subseteq \overline{\a_{pt}} \text{ for } t \gg 1\right\}
	\ \in\ \RR_{\ge0} \cup \{\infty\},
	\]
	with the convention that $\inf \emptyset = \infty$. For ideals $\a,\b\subseteq R$ we write $\Loj_{\b}(\a)$ for $\Loj_{\{\b^q\}_{q \ge 1}}\bigl(\{\a^p\}_{p\ge1}\bigr)$.
\end{definition}

\begin{theorem}[Valuative lower bound]\label{thm:valuative-filtrations}
	Assume that $(R,\m)$ is a Noetherian local domain. Let $\abul=\{\a_p\}_{p\ge1}$ and
	$\bbul=\{\b_q\}_{q\ge1}$ be graded families of ideals in $R$.
	Then
	\[
	\Loj_{\bbul}(\abul)\ \ge\ \sup_{v}\ \frac{v(\abul)}{v(\bbul)},
	\]
	where the supremum is taken over all valuations $v$ centered at $\m$ with $v(\bbul) > 0$.
\end{theorem}

\begin{theorem}[Finite Rees valuations for a Noetherian filtration]
	\label{thm:finite-testing-fingen}
	Let $(R,\m)$ be an excellent Noetherian local domain, and let
	$\abul=\{\a_p\}_{p\ge1}$ be a filtration of $\m$-primary ideals in $R$.
	Assume that the Rees algebra $\R(\abul):=\bigoplus_{p\ge0}\a_pT^p$
	is finitely generated as an $R$--algebra.
	Let $\overline{\R(\abul)}$ denote its integral closure, and set $\overline{Y}:=\Proj \overline{\R(\abul)}.$
	Then, there exists a finite set of divisorial valuations
	\[
	\mathcal V(\abul)=\{v_1,\dots,v_r\},
	\]
	arising from prime divisors on $\overline{Y}$ lying over $\m$, such that for every filtration $\bbul=\{\b_q\}_{q\ge1}$ of ideals in $R$ satisfying $v_i(\bbul) > 0$ for all $1 \le i \le r$, one has
	\[
	\Loj_{\bbul}(\abul)
	=
	\max_{1\le i\le r}
	\frac{v_i(\abul)}{v_i(\bbul)}.
	\]
\end{theorem}

\begin{proof}
	We break the proof into two steps.
	
	\smallskip
	\noindent\textit{Step 1: a finite divisorial testing set for the integral closures
		$\overline{\a_p}$.}
	Let
	\[
	\R:=\R(\abul)=\bigoplus_{p\ge 0}\a_pT^p
	\subseteq R[T]
	\quad\text{and}\quad
	\overline{\R}:=\overline{\R(\abul)}
	\]
	be the integral closure of $\R$ in its total ring of fractions.
	Since $\R$ is finitely generated over the excellent Noetherian domain $R$, it is Noetherian,
	and $\overline{\R}$ is a finite $\R$--module; in particular
	$\overline{\R}$ is again a finitely generated $R$--algebra.
	
	Let $\pi:\overline{Y}\to \Spec R$ the induced proper birational morphism.
	Write $E\subseteq \overline{Y}$ for the exceptional locus of $\pi$.
	Since $\overline{Y}$ is of finite type over $R$, the exceptional locus has only finitely many
	irreducible components of codimension one; denote them by
	\[
	E_1,\dots,E_r.
	\]
	Each $E_i$ defines a divisorial valuation $v_i$ of $\Frac(R)$ by
	\[
	v_i(h):=\ord_{E_i}(\pi^*h)\quad (0\ne h\in \Frac(R)),
	\]
	normalized so that $v_i(R)\ge 0$ and $v_i(\m)>0$ whenever $\pi(E_i)\subseteq V(\m)$.
	
	We claim that these finitely many valuations test the integral closure of each
	$\a_p$ in the sense that
	\begin{equation}\label{eq:ic-test-by-Ei}
		\overline{\a_p}
		=
		\{\,h\in R:\ v_i(h)\ge v_i(\a_p)\ \text{for all } i=1,\dots,r\,\}.
	\end{equation}
	Indeed, for $h\in R$ we have
	\[
	h\in \overline{\a_p}
	\ \Longleftrightarrow\
	hT^p \in \overline{\a_pT^p}\subseteq \overline{\R}_p
	\ \Longleftrightarrow\
	hT^p \in \overline{\R},
	\]
	where $\overline{\R}_p$ denotes the degree-$p$ piece of $\overline{\R}$.
	Now, $\overline{\R}$ is a normal Noetherian domain, so a Krull domain; therefore,
	\[
	\overline{\R}
	=
	\bigcap_{\substack{\mathfrak P\in \Spec \overline{\R}\\ \mathrm{ht}(\mathfrak P)=1}}
	(\overline{\R})_{\mathfrak P}
	\ \subseteq\ \Frac(\overline{\R}).
	\]
	Consequently, $hT^p\in \overline{\R}$ if and only if $hT^p$ lies in each DVR
	$(\overline{\R})_{\mathfrak P}$ for $\mathrm{ht}(\mathfrak P)=1$.
	Among the height-one primes $\mathfrak P$, only those corresponding to the codimension-one
	components of the exceptional locus contribute nontrivially to membership in $\a_p$; the
	corresponding valuations are precisely $v_1,\dots,v_r$, and the resulting condition is
	exactly \eqref{eq:ic-test-by-Ei}.  This proves the claim.
	
	We set
	\[
	\mathcal V(\abul):=\{v_1,\dots,v_r\}.
	\]
		
		\smallskip
		\noindent\textit{Step 2: computation of $\Loj_{\bbul}(\abul)$.}
		Fix any filtration $\bbul=\{\b_q\}_{q\ge 1}$ of ideals in $R$.
		Set
		\[
		\lambda\ :=\ \max_{1\le i\le r}\ \frac{v_i(\abul)}{v_i(\bbul)}\ \in\ \RR_{\ge 0}.
		\]
	
		Let $\frac{q}{p}\in\QQ_{>0}$ be admissible in Definition~\ref{def:L-families}, i.e.,
		\[
		\b_{qt}\subseteq \overline{\a_{pt}}
		\quad\text{for all integers }t\gg 1.
		\]
		By \eqref{eq:ic-test-by-Ei} applied to $pt$, this containment is equivalent to
		\[
		v_i(\b_{qt})\ \ge\ v_i(\a_{pt})
		\quad\text{for all }i=1,\dots,r\ \text{and all }t\gg 1.
		\]
		By the same arguments as in the proof of \Cref{thm:valuative-filtrations}, we get
	$q\,v_i(\bbul)\ge p\,v_i(\abul)$, i.e.,
		\[
		\frac{q}{p}\ \ge\ \frac{v_i(\abul)}{v_i(\bbul)}\quad\text{for all } i.
		\]
	Hence, $\frac{q}{p}\ge \lambda$. Taking the infimum over admissible \(\frac{q}{p}\) yields
		\(\Loj_{\bbul}(\abul)\ge \lambda\).
		
		Conversely, fix $\varepsilon>0$ sufficiently small. 
		By the definition of $v_i(\abul)$ and $v_i(\bbul)$, there exists $p_0$
		such that for all $p\ge p_0$ and all $i$,
		\[
		\frac{v_i(\a_p)}{p}\ \le\ v_i(\abul)+\varepsilon
	\text{ and }
		\frac{v_i(\b_p)}{p}\ \ge\ v_i(\bbul)-\varepsilon > 0.
		\]
    Set $$\delta = \epsilon \cdot \max\limits_{1 \le i \le r}\frac{v_i(\abul) + v_i(\bbul)}{v_i(\bbul)(v_i(\bbul)-\epsilon)}.$$
    Then, $\delta \rightarrow 0$ as $\epsilon \rightarrow 0$ and 
    $(\lambda + \delta)(v(\bbul)-\epsilon) \ge v_i(\abul)+\epsilon.$
		Choose $p\ge p_0$ large and set
		$$q:=\left\lceil(\lambda+\delta)p\right\rceil.$$
		Now choose $t_0$ large enough so that $qt\ge p_0$ for all $i$ and all $t\ge t_0$.
		Then, for every $i$ and every integer $t\ge t_0$ we have
		\[
		v_i(\b_{qt})\ \ge\ qt\,(v_i(\bbul)-\varepsilon)
		\ \ge\ (\lambda+\delta)pt\,(v_i(\bbul)-\varepsilon)
		\ \ge\ pt\,(v_i(\abul)+\varepsilon)
		\ \ge\ v_i(\a_{pt}).
		\]
		Hence, by \eqref{eq:ic-test-by-Ei},
		\[
		\b_{qt}\subseteq \overline{\a_{pt}}
		\quad\text{for all integers }t\ge t_0.
		\]
		Thus, $\frac{q}{p}$ is admissible in Definition~\ref{def:L-families}, so
		\[
		\Loj_{\bbul}(\abul)\ \le\ \frac{q}{p}\ \le\ \lambda+\delta+\frac{1}{p}.
		\]
		Letting $p\to\infty$ and then $\varepsilon\to 0$ ($\delta \to 0$) gives $\Loj_{\bbul}(\abul)\le \lambda$.
		The theorem is proved.
	\end{proof}

\end{document}
